\section*{Chapter \ref{ch:qm:brownian-motion} --- Brownian Motion}

\begin{solution}{exo:qm:brownian-motion:1}
$\Cov(W_2, W_5) = \min(2,5) = 2$; $\Var(W_5 - W_2) = 3$; correlation $2/\sqrt{2 \times 5} = 0.632$.
\end{solution}

\begin{solution}{exo:qm:brownian-motion:2}
By the reflection principle, $\P(M_1 \ge 1) = 2(1 - \Phi(1)) = 31.7\%$.
\end{solution}

\begin{solution}{exo:qm:brownian-motion:3}
$0.20/\sqrt{252} = 1.26\%$ a day; $\sqrt5$ times that, 2.82\%, a week.
\end{solution}

\begin{solution}{exo:qm:brownian-motion:4}
$b = \ln 0.95 = -0.0513$, $T = 0.25$: continuously $2\Phi(b/0.15) = 73.2\%$; at the closes, the
level shifted by $0.0110$ gives 67.8\%.
\end{solution}

\begin{solution}{exo:qm:brownian-motion:5}
$\E_s[W_t^2] = W_s^2 + (t - s)$, so $W_t^2 - t$ is a \omterm{def:qm:probability-at-speed:martingale}{martingale}. With $\tau$ the exit time,
$\E[W_{\tau\wedge n}^2] = \E[\tau \wedge n]$; letting $n \to \infty$ (bounded convergence on the
left, monotone on the right) and using $\P(W_\tau = b) = a/(a+b)$: $\E[\tau] = b^2\frac{a}{a+b} +
a^2\frac{b}{a+b} = ab$.
\end{solution}

\begin{solution}{exo:qm:brownian-motion:6}
$\mathcal N(0, \tfrac14)$, from the covariance $s(1 - s)$; by
\cref{prop:qm:brownian-motion:between}, $\P(\max \ge 0.5) = \exp(-2 \times 0.5 \times 0.5) = 60.7\%$.
\end{solution}

\begin{solution}{exo:qm:brownian-motion:7}
$\Var(\sum\Delta_k^2) = \sum 2(1/n)^2 = 2/n$: at $n = 4\,096$ the standard deviation is 0.022, and
2\,000 simulated paths give 0.022. The figure's value at 4\,096 points, 1.031, is 1.4 standard
deviations from 1: consistent.
\end{solution}

\begin{solution}{exo:qm:brownian-motion:8}
Checking daily closes prices a contract monitored daily, not continuously: it misses every
crossing between closes and underestimates the knock-out probability by several points (72\%
against 82\% for the 2\% stop of the chapter). More paths only reduce the statistical error,
not this bias, which falls like $\sqrt{\Delta t}$. Add each interval's bridge crossing
probability, or price the continuous contract with the level shifted by $\beta\sigma\sqrt{\Delta t}$
to check.
\end{solution}

\begin{solution}{pb:qm:brownian-motion:1}
\textbf{1.} $b = \ln 0.98 = -0.0202$; $\sigma\sqrt T = 0.30\sqrt{21/252} = 0.0866$.
\textbf{2.} $\Phi(-0.0202/0.0866) = \Phi(-0.233) = 40.8\%$.
\textbf{3.} $2 \times 40.8\% = 81.6\%$.
\textbf{4.} By the reflection principle each path that touches the stop and ends above it is
paired with its mirror image, which ends below.
\textbf{5.} 82.4\%: the negative drift makes the touch slightly likelier.
\textbf{6.} 72.0\% (standard error 0.1 point).
\textbf{7.} $0.5826 \times 0.30/\sqrt{252} = 0.0110$ in log price: the stop behaves like a 3.1\%
stop watched continuously.
\textbf{8.} 71.9\%.
\textbf{9.} 79.6\% by simulation, 79.6\% by the shift.
\textbf{10.} The overshoot, and hence the error, is of order $\sigma\sqrt{\Delta t}$: dividing
$\Delta t$ by four halves it (77.1\%, 79.1\%, 80.5\% at 4, 16 and 64 observations a day).
\textbf{11.} $\exp(-2 \times 0.01 \times 0.01/(0.09/252)) = 57.1\%$.
\textbf{12.} 81.6\% (standard error 0.08 point), equal to the continuous answer.
\textbf{13.} Given the closes, the crossings in different days are independent bridges, so $1
- \prod(1 - p_k)$ is the conditional probability of a touch; its mean is the unconditional one.
\textbf{14.} The bridge: it is exact at the cost of daily steps, while 64 steps a day still
leave a 1-point bias at 64 times the cost.
\textbf{15.} 64.1\%.
\textbf{16.} 89.3\% continuously, 83.5\% at the closes.
\textbf{17.} The continuous number, 82\%; the model misses gaps (the order fills below the stop
when the price jumps through it), volatility that is not constant, and intraday patterns.
\textbf{18.} The daily-close number, 72\%: the contract is monitored at the closes.
\textbf{19.} Named result: \emph{the stop-loss}: a 2\% stop on a 30\%-volatility stock is touched
within a month with probability 81.6\% if watched continuously and 72.0\% if watched at the
closes; the Broadie--Glasserman--Kou shift gives 71.9\% and the bridge-corrected daily
simulation 81.6\%.
\textbf{20.} Because a Brownian path crosses and re-crosses a level between observations, and
the closes see only the crossings that last until the close.
\end{solution}

\begin{solution}{iq:qm:brownian-motion:1}
$W$ is a bounded \omterm{def:qm:probability-at-speed:martingale}{martingale} up to the exit, so $\E[W_\tau] = 0$: $-1 \times p + 2(1 - p) = 0$,
$p = 2/3$.
\iqlookfor{optional stopping in one line.}
\end{solution}

\begin{solution}{iq:qm:brownian-motion:2}
The squared increments over a partition have mean equal to the elapsed time and a variance
that vanishes with the mesh. It is why the variance of a hedged book accumulates with time
whatever the rebalancing frequency, why volatility is measurable from one path, and why
Itô's formula has a second-order term.
\iqlookfor{the $L^2$ computation and its consequences.}
\end{solution}

\begin{solution}{iq:qm:brownian-motion:3}
By the reflection principle $M_1$ has the law of $|W_1|$, so $\E[M_1] = \sqrt{2/\pi} = 0.798$.
\iqlookfor{reflection, then the half-normal mean.}
\end{solution}

\begin{solution}{iq:qm:brownian-motion:4}
Simulate on the grid you need anyway and, in each interval, multiply by the bridge
probability of not crossing, $1 - \exp(-2(x_k - b)(x_{k+1} - b)/\sigma^2\Delta t)$ (or sample the
crossing with that probability); for a quick check, shift the level by $0.5826\,\sigma\sqrt{\Delta t}$.
\iqlookfor{the Brownian-bridge crossing probability.}
\end{solution}

\begin{solution}{iq:qm:brownian-motion:5}
The probability of ending the month below the stop, 41\%. A stop is triggered by the minimum,
not the terminal value; by reflection the right number is twice that, 82\% continuously, 72\%
if only closes count.
\iqlookfor{terminal versus path-dependent probability.}
\end{solution}

\begin{solution}{iq:qm:brownian-motion:6}
$(W_1, W_2 - W_1)$ are independent standard normals $(X, Y)$, and the event is $X > 0$, $X + Y >
0$: a wedge of angle $\pi/2 + \pi/4$ in the plane, so the probability is $(3\pi/4)/(2\pi) = 3/8$.
\iqlookfor{rotational symmetry of the Gaussian.}
\end{solution}
